\documentclass[11pt]{article}
\usepackage[T1]{fontenc}
\usepackage{lmodern,microtype,amsmath,amssymb,amsthm,booktabs}
\usepackage[margin=1in]{geometry}
\usepackage[hidelinks]{hyperref}
\newtheorem{lemma}{Lemma}
\newtheorem{proposition}{Proposition}
\newtheorem{theorem}{Conditional theorem}
\newcommand{\FF}{\mathbb F_2}
\newcommand{\QQ}{\mathbb Q}
\title{Geometric envelope frames for a smaller complex network}
\author{Research exploration prepared with OpenAI Codex}
\date{October 8, 2026}
\begin{document}
\maketitle
\begin{abstract}
We propose replacing the complex motif's separate even-intersection side
wires by two cancellation-free sum circuits. Coordinate envelopes, rather
than spans of all input indicators, label the disjointness circuit. A single
unused coordinate guarantees nonalternating residuals at its output boundary.
For 25 ground points and a 26-dimensional binary label space, the exact finite
construction supports a complex saving $4\cdot10^{-9}$. Combined with the
retained bit network and compact-control estimates, the supplied parameter
calculation gives $\kappa=5.9\cdot10^{-10}>2^{-31}$, about $7.11$ times the
repository's $8.3\cdot10^{-11}$. This is an unreviewed conditional research
argument, not a formally verified multiplication theorem or a runtime result.
\end{abstract}

\section{Scope, attribution, and what is changed}
The baseline is CrocSwap/integer-mult-bounds at commit
\begin{center}\texttt{6e564879f51ae16f23d392e9e196c605f36d90df}.\end{center}
It pins the OpenAI manuscript at
\begin{center}\texttt{adc7f1241b42e322a6451854ab7e4b4c146bf78a}.\end{center}
We retain its fixed finite-alphabet, fixed-number-of-one-dimensional-tapes
model; the paired bit exponent $\tau=1-296/10^{11}$; nonadjacent routing;
compact-control movement, reservation, repair, and recurrence arguments;
tight Gaussian setup; and all downstream analytic and rounding interfaces.
These retained results remain assumptions here.

Only the finite complex network is replaced. The modifications are
(i) shared scalar sum circuits, (ii) geometric envelope labels, and
(iii) one extra \emph{label-space} coordinate. The extra coordinate increases
the motif arity from $25^3$ to $26^3$. It is not a free address coordinate,
scratch tape, or clean register. The layer groups its already existing axes
into this new arity, with its usual padding and complete-row accounting.
Every scalar auxiliary value is arbitrary and restored.

Restored auxiliary circuits and complement frames are adapted from the
repository's incidence, DAG, and paired-bit constructions. Monotone
disjoint-set summation is established prior work; see Kaski, Koivisto, and
Korhonen, \emph{Fast Monotone Summation over Disjoint Sets},
\href{https://arxiv.org/abs/1208.0554}{arXiv:1208.0554}.
Our finite experiment uses the simple exact circuit defined below, not that
paper's asymptotic circuit-size bound. No claim of global novelty or optimality
is made. The new point of this exploration is the binary residual argument
for this particular complex interface.

\section{Two exact scalar circuits}
Let $n=25$, $\mathcal T=\binom{[n]}3$, $v=2300$, $h=n+1=26$, and
$F=\FF^h$ with its ordinary dot product. Triple indicators $t_T$ have zero
in the extra coordinate, denoted $\star$. Define
\[
 (D x)_S=\sum_{T\cap S=\varnothing}x_T,\qquad
 (E x)_S=\sum_{|T\cap S|=2}x_T.
\]
The central gather/scatter is retained on the $n$ ground points and its
constant channel: its coefficient from $T$ to $S$ is
$(|T\cap S|-1)/2$. Therefore
\[
 \frac12D-\frac12E+R_0G=I.
\]
This is an identity over $\mathbb Z[i,1/2]$: the four intersection sizes
$0,1,2,3$ give respectively $0,0,0,1$ on the left.

For $D$, the output supports are the families of triples disjoint from each
target. For $E$, fix a pair $P$ and, for each target triple $S\supset P$,
form the partial output consisting of all triples containing $P$ except $S$.
Inject these partial outputs into $S$. Each intersection-two pair occurs
exactly once, indexed by its common pair.

Here is a complete finite rule for constructing both DAGs. Order all triples
lexicographically and assign them indices $0,\ldots,v-1$. A singleton support
is its input variable. For any larger requested support, take the smallest
dyadic interval of indices containing it and split at that interval's midpoint.
Recursively construct the two nonempty parts and add them. Intern equal
supports, and discard nodes not reaching any requested output. Every addition
combines disjoint formal supports, so every coefficient is exactly one.
This description is deterministic and independent of numerical test inputs.

The companion program checks every output support, giving
\begin{center}
\begin{tabular}{lrrr}
\toprule
Circuit & Additions & Partial outputs & Reversible roles\\
\midrule
$D$ & 212737 & 2300 & 215037\\
$E$ & 36620 & 6900 & 43520\\
\bottomrule
\end{tabular}
\end{center}
Thus the total number of side roles per invocation is $R=258557$.
The baseline separate-pair construction on the same 25 points requires
$v(\binom{22}3+66)=3693800$ side roles. The new label dimension is larger;
all subsequent counts use $h=26$ rather than silently retaining $25$.

\section{Reversible embedding and arbitrary scratch}
For a DAG with $c$ binary additions and $q$ designated output occurrences,
allocate one role per outgoing use. At each addition identify one incoming
role with one outgoing pivot; the other incoming role retires. At an input,
copy its pivot into its other outgoing roles. At an addition, add its second
incoming role into its pivot, then copy the pivot into its other outgoing
roles. The unused destinations may be dirty: each gate is an invertible
addition. There are $2c+q$ uses and $c$ identifications, hence $c+q$ roles.
Different simultaneously active roles are never identified.

Write $L$ for the direct sum of the two reversible DAG programs, $V$ for
source loading, and $J$ for injection with weights $+1/2$ on $D$ and $-1/2$
on $E$. The chronological forward schedule is
\[
 L,\ -J,\ L^{-1},\ -R_0,\ V,\ G,\ R_0,\ L,\ J,\ L^{-1},\ -G,\ -V.
\]
Here a minus on an additive update denotes its inverse. For initial side
scratch $z$, the first injection contributes $-JLz$, and the second
contributes $JL(z+Vx)$. Cleanup restores $z$. The central scratch cancels
in the same way. The net target change is
$(JLV+R_0G)x=x$, with sources and all scratch restored. The inverse schedule
gives $y\gets y-x$. Three tensor stages, forward/inverse/forward with the
logical banks interchanged in the middle, give $(X,Y)\mapsto(-Y,X)$.
Use distinct scratch banks for all invocations; no stage-sharing assumption
is needed for the counts below.

\section{The geometric residual lemma}
Merely copying the rational span argument is invalid over $\FF$: a subspace
may be degenerate, or a nonzero orthogonal residual may be alternating.
The upstream complex phase decomposition needs an orthonormal basis for
every nonzero residual.

\begin{lemma}[Envelope labels]
For each node $z$ of $D$, use $U_z=\langle t_T\rangle$ if it is a singleton
input, and otherwise use the coordinate space on the union of its source
triples. For a node of $E$, use the span of its source triple indicators.
These labels are nondegenerate, are nested along each DAG edge, and lie in
$t_S^\perp$ for every output injected into target $S$. Every nonzero residual
on a side-role path, including its entry, output, and cleanup, is
nonalternating. The same holds for reverse paths labeled by $U_z^\perp$.
\end{lemma}
\begin{proof}
A triple line has norm one. All non-input $D$ labels are coordinate spaces,
and their supports have at least four coordinates, since they include two
distinct triples. Coordinate spaces are nondegenerate and nested whenever
their supports are nested. A line $\langle t_T\rangle$ entering such a space
has a complement containing a coordinate unit outside $T$, so that complement
is nonalternating. Between two coordinate spaces the residual is another
coordinate space. A $D$ output uses coordinates disjoint from its target.

All triples contributing to a nonsingleton $E$ node contain a fixed pair.
Distinct triples in that pair-star have intersection two, so their indicators
form an orthonormal family. Any subset spans a nondegenerate space; every
difference between nested such spans is spanned by a subset of the same
orthonormal family. A support shared between pair-stars is either a singleton
or has the same unique common pair, so interning does not invalidate this
argument. Output source indicators are orthogonal to their target indicator.

In either circuit, $U_z$ and $\langle t_S\rangle$ are nondegenerate and
orthogonal at an output. Thus
$t_S^\perp\cap U_z^\perp$ is nondegenerate. It contains the unit vector
$e_\star$, since no source or target indicator and no coordinate envelope
uses $\star$. It is therefore nonalternating. The same unit vector lies
in every cleanup complement $U_z^\perp$. Entry labels themselves have
orthonormal bases as just described. Finally, complementing labels and
reversing an inclusion preserves its orthogonal residual. This proves the
reverse-direction statement.
\end{proof}

The extra coordinate has a specific, testable purpose. Without it, a full
disjointness output has envelope equal to the coordinates outside $S$.
Its residual inside $t_S^\perp$ is the even-weight plane on the three
coordinates of $S$, a nonzero alternating space. It has no norm-one vector.
Our verifier explicitly rejects this tempting but invalid construction.

\section{Tensor-stage frames and accounting}
Use the original three tensor stages, now in $\mathcal F=F^{\otimes3}$,
with $m=h^3=17576$. At stage $j$, write
$A=F^{\otimes(j-1)}=B\perp P$, where $P$ is the earlier tensor indicator
line, and let $Q$ be the future indicator line. Suppress the factor $Q$ and
write
\[
 D_0=B\otimes F,\quad D_1=A\otimes F,\quad
 X_t=D_0\perp(P\otimes\langle t\rangle),\quad
 Y_t=D_0\perp(P\otimes t^\perp),\quad
 D_U=D_0\perp(P\otimes U).
\]
For the forward schedule in the preceding section the labels are, in order,
\[
 D_0,D_0,D_0,D_0,X_t,D_1,D_0,D_{U_z},Y_t,D_1,D_1,D_1.
\]
The middle mixer has a separate label for each DAG gate. Source and target
updates are grouped by physical triple, with the same frame on all incidences.
New mixer roles enter at $D_0$; retired ones next move to $D_1$ in cleanup.
The lemma proves all side changes increase with admissible residuals.

For the inverse schedule with exchanged logical banks the chronological
operations are
\[
 V,G,L,-J,L^{-1},-R_0,-G,-V,R_0,L,J,L^{-1},
\]
with labels
\[
 D_0,D_0,D_0,X_t,D_{U_z^\perp},D_1,D_0,Y_t,D_1,D_1,D_1,D_1.
\]
Here the middle inverse mixer visits nodes in reverse topological order.
An output slot enters it from its physical $X$ target line, contained in
$U_z^\perp$. An input slot exits at its physical $Y$ source complement.
This checks the often-missed reverse frame obligation.

The physical data boundaries are unchanged in form:
$X$ goes from $A\otimes\langle t\rangle$ to $A\otimes F$ and $Y$ from
$B\otimes\langle t\rangle$ to $Y_t$, all tensored with $Q$.
The boundary changes involving $B$, a triple complement, or a future
complement have the original nondegenerate tensor decompositions.
Whenever $B$ or $Q^\perp$ is nonzero, a coordinate outside the respective
tensor indicator support gives a norm-one vector. Tensoring these vectors
with the norm-one vectors from the lemma proves that every nonzero residual
in the full network is nonalternating. A nondegenerate nonalternating binary
symmetric space has an orthonormal basis, as proved in the upstream motif
section.

Only a central return $D_1\to D_0$ decreases. There are $n+1=h$ center roles
per invocation, each losing dimension $h$, and $I=3v^2$ invocations. Hence
\begin{align*}
 N&=v^3=12167000000,\\
 W&=2N+3v^2(R+n+1)=4128046210000,\\
 L_{\mathrm{dec}}&=3v^2(n+1)h=10728120000,\\
 s&=Wm-2N+2L_{\mathrm{dec}}=72554537309200000.
\end{align*}
This counts all physical roles and all edge residuals, including cleanup;
the source and sink dimension sums telescope exactly as in the baseline.

\begin{proposition}[Proposed complex interface]
The new finite network has the original complex phase interface with the
above $m,W,s$. Its relative deficit is
\[
 \eta=\frac{Wm-s}{Wm}=\frac{68}{1714426753}>0,
 \qquad \frac{s}{W}<m^{1-4\cdot10^{-9}}.
\]
\end{proposition}
\begin{proof}
Scalar restoration, both frame orientations, and every residual type were
checked above. For a nested edge, orthogonal weight additivity modulo four
and an orthonormal residual basis express its phase change as one $C$ or
$C^{-1}$ translation kernel per basis vector, exactly as in the original
interface. All triple tensors still have Hamming weight $27$, so the original
source/sink sign corrections and signed bank exchange apply to every data
and auxiliary role. The extra coordinate changes ambient dimension but
does not change these endpoint weights. Exact rational logarithm bounds
certify $\eta>(4\cdot10^{-9})\log m$. Since $e^{-x}>1-x$ for $x>0$,
the strict recurrence comparison follows.
\end{proof}

\section{Coefficient guard and completed layer}
The old assertion that each wire touches at most twelve scalar gates does
not apply to a reused DAG pivot. We supply a replacement operation count.
One mixer uses at most $2R$ elementary additions, so its four executions
cost at most $8R$. Two source loads use at most $4v$ additions; two weighted
injections cost at most $4q$ additions/halvings, with $q=4v$. Two central
gathers and two scatters cost at most $18v$. Thus the deliberately loose
bound $32(R+v+n+1)$ suffices per invocation. The verifier checks
\[
 32I(R+v+n+1)+4s+4W+4 < E_0:=64(W+m+1)^3.
\]
Every term is charged per coefficient; permutations do not add arithmetic
depth. Also $2\le s<m^5$. Consequently the retained guard recurrence
$A(e)\le sA(e/m)+E_0$, with $A(e)\le8e$ at a leaf, applies with these new
constants. In particular $C_1=5-4\beta+\zeta$ is still valid. Recompute
$B=s+E_0$ and
\[
 C_0=\left\lceil\max\{128mB^2,18mB^2(1+1/\zeta)\}\right\rceil.
\]
These values and the guard inequalities are in the exact certificate.

No change to compact controls is required: this network supplies a fixed
finite collection of the same admissible translation kernels, and the bit
network implements the required binary basis changes. Distinct bit and
complex arities are already allowed. Use
\begin{gather*}
 \tau=1-\frac{296}{10^{11}},\qquad \sigma=1-\frac4{10^9},\\
 \epsilon=\frac{1999}{10000},\quad c=1,\quad
 \beta=\frac1{1000},\quad\zeta=\frac1{10000},\quad\delta=\frac1{10^6},\\
 C_1=\frac{49961}{10000},\quad
 \lambda=1-\frac{2959}{10^{12}},\quad
 \lambda'=1-\frac{2958}{10^{12}},\quad
 \kappa=\frac{59}{10^{11}}.
\end{gather*}
The compact-control recurrence allows either exponent ordering. Here
$\sigma<\tau$, so its internal exponent is $\chi=\tau$, its leaf exponent
is $1-3996/10^{12}$, and its reservation exponent is $\max(1-c,0)=0$.
All are below the required intermediate exponents with strict gaps.
At $c=1$, reservation preprocessing is $O(V(\log d+\log p+1))$.
The written reservation proof permits any fixed $c>0$; it does not require
$c<1$. The other address conditions remain valid since
$\epsilon(1+c)=3998/10000<1$, and
$\epsilon C_1=99872039/10^8<1$.

The seven retained assembly margins are
\begin{align*}
 g_1&=1-\epsilon(1+c),&g_2&=\epsilon c(1-\tau),\\
 g_3&=\epsilon(1-\lambda'),&g_4&=(1-\tau)(1-\epsilon),\\
 g_5&=1/4-\delta-5\epsilon/4,&g_6&=1-\delta-\epsilon,&g_7&=\epsilon.
\end{align*}
Their exact minimum is
\[
 G_*=g_3=\frac{2956521}{5\cdot10^{15}}=5.913042\cdot10^{-10},
 \qquad G_*-\kappa=1.3042\cdot10^{-12}>0.
\]
All other setup inequalities are also checked with rational arithmetic.

\begin{theorem}
Conditional on the retained upstream and compact-control interfaces and the
finite-network argument above, the retained assembly yields
\[
 T(\ell)=O\!\left(\ell(\log\ell)^{1-59/10^{11}}\right).
\]
In particular the dyadic corollary $\kappa=2^{-31}$ follows.
\end{theorem}
\begin{proof}
The new motif and recomputed guard supply the same completed-layer contract.
Its internal, leaf, reservation, repair, and padding costs are covered by the
retained compact-control estimate. The setup inequalities and seven strict
assembly margins therefore give the conclusion; their strict positive gap
absorbs the fixed logarithmic factors. No downstream analytic estimate has
been strengthened in this inference.
\end{proof}

The saving ratio is $590/83\simeq7.10843$. This measures an asymptotic exponent,
not practical speed. Once this complex bottleneck is removed, the retained
bit saving and $\epsilon<1/5$ imply the scoped ceiling
$\kappa<(296/10^{11})/5=5.92\cdot10^{-10}$. The proposed witness reaches
over $99.6\%$ of that ceiling. Further substantial improvement needs another
bit network, layer construction, or assembly argument.

\section{Audit of the complete invocation and phases}
The side-DAG inclusions alone do not establish the full network score. The
additional program \path{audit.py} follows every physical data, center, and
side role through the complete twelve-step schedule, its reverse, and all
three tensor stages at ground sizes six and seven. It constructs residual
orthonormal bases, including the absorption of alternating planes, and checks
at every binary address that the actual frame phase equals the product of
its proposed $C_v$ or $C_v^{-1}$ phases. The comparison uses exact integers
modulo four. This tests the phase factorization itself, not only ranks or
nonalternation.

If $a=h^{j-1}$, the two stage components are $B\otimes F$ and $P\otimes F$.
The program follows their projectors separately, multiplying the $B$ rank
by $a-1$. It also charges the full future complement at each scratch sink.
The resulting physical-edge sum is independently compared with
\[
 (R+h)m+2va(h-1)+2h^2
\]
for every orientation and stage. The only decreases have total $h^2$.
Multiplying the data contributions by $v^2$ and summing the three stages
recovers $Wm-2N+2L_{\mathrm{dec}}$. Thus the rank certificate has a separate
check on the whole schedule and its tensor accounting.

A further bounded experiment varied the first/third and middle bit-network
sizes independently. Keeping equal outer sizes allows the existing matched
auxiliary-bank sharing. For outer size $x$ and middle size $y$, put
$v=\binom{x}3$, $u=\binom{y}3$; the proposed counts in this experiment are
\[
 N=v^2u,\quad m=x^2y,\quad
 L=2vux^2+v^2y^2,
\]
\[
 W=2N+vu(R_x+x)+v^2(R_y+y).
\]
The script \path{bit_screen.py} computes exact logarithmic enclosures for
all 441 choices with $x,y\in\{30,32,\ldots,70\}$. The retained $(50,50,50)$
network is the unique winner in this finite screen. Its full circuit and both
frame orientations are checked again. This experiment supplies no new bit
saving and asserts no all-size or all-network optimality theorem.

\section{Reproduction and proof boundary}
Run Python 3.11 or newer from the repository root:
\begin{verbatim}
python3.11 research/geometric-complex/verify.py
\end{verbatim}
The verifier checks all source supports and scalar coefficients of both
full-size side DAGs; every compiled role incidence in both directions;
the geometric hypotheses for every full-size node; direct binary projector
residuals at small sizes; every scalar input-basis direction, including dirty
scratch, at sizes six and seven; signed three-shear composition; the invalid
missing-coordinate control; the new guard; and all assembly margins.
The JSON includes source hashes, retained-dependency hashes, circuit hashes,
and the complete physical-path/phase audit. The stronger witness
$59/10^{11}$ replaces the initial exploratory $58/10^{11}$ after checking
all strict assembly gaps. Run the additional bounded bit experiment with
\begin{verbatim}
python3.11 research/geometric-complex/bit_screen.py
\end{verbatim}

Finite evidence does not by itself prove the general tape theorem.
The tensor residual argument, precision integration, and application of the
retained completed-layer interface above are written mathematical arguments
requiring independent review. No full multiplication-machine implementation,
formal proof, or combined manuscript patch is supplied. The repository's
headline, pinned upstream source, and historical certificates are unchanged.
This work should be treated as a reviewable conditional candidate, not an
established new record.
\end{document}
